% Section II: literature review (2022--2026). Included in main.tex after the introduction.
\section{Literature Review}
\label{sec:lit}

\textbf{Malicious-package detection.} Most supply-chain defences judge each published version on its own. Attack
catalogues and taxonomies document account takeover, typosquatting and install-script abuse~\cite{ohm2020backstabber,
ladisa2023sok}, and ecosystem studies show how a few maintainer accounts can reach much of npm~\cite{zimmermann2019small,
zahan2022weak}. Amalfi combines classifiers, a check that a package rebuilds from its source repository, and clone
detection, and found 95 previously unknown malicious npm versions in one week of publications~\cite{sejfia2022amalfi}.
Ladisa et al.\ train one classifier on language-independent features for npm and PyPI and report 58 new malicious
packages~\cite{ladisa2023crosslang}. These scanners run on the registry side and ask whether a version looks malicious;
they do not compare a version with the one a user already trusts.

\textbf{Signing and provenance.} Sigstore~\cite{newman2022sigstore}, SLSA~\cite{slsa2023provenance}, npm
provenance~\cite{npmprovenance} and PEP~740~\cite{pep740} let a consumer check who built a release and from which
repository. Adoption depends on registry policy: Schorlemmer et al.\ find that 0.2\% of PyPI artifacts were signed in
2023, before PEP~740 attestations, and that mandates raise the number of signatures while dedicated tooling raises their
quality~\cite{schorlemmer2024signing}. Provenance tells who built a release; it does not say whether the builder changed
since the last version a user accepted. MCP-Lock uses build identity for that continuity check, not as proof that a
release is safe.

\textbf{Lockfiles.} Gamage et al.\ compare the lockfiles of seven package managers. Most record versions and checksums,
and npm and Cargo also record a source link~\cite{gamage2025lockfiles}; publisher identity, install scripts and the
tools a package advertises are not among the contents they describe. An MCP configuration that launches \code{npx -y pkg} has no lockfile at
all.

\textbf{MCP security.} Surveys and taxonomies describe tool poisoning, rug pulls and malicious
servers~\cite{hou2025mcp,zhao2025attack,gaire2025sok}, and practitioners described the rug-pull pattern
early~\cite{invariant2025poisoning}. Gaire et al.\ trace rug pulls to mutable server logic, the absence of continuous
integrity checks and the absence of re-approval triggers, and discuss signed, immutable tool
manifests~\cite{gaire2025sok}. Huang et al.\ test tool poisoning on seven MCP clients and find that five do not
validate tool metadata statically; they recommend signed, immutable tool definitions and checks at
registration~\cite{huang2026toolpoison}. Hasan et al.\ study the security and maintainability of open-source
servers~\cite{hasan2025glance}. Kraishan measures drift of the fields the MCP registry versions across 8,900
multi-version servers, notes that tool definitions are not versioned there, and recommends client-side
pinning~\cite{kraishan2026drift}; we measure the tool definitions themselves and build that pinning. Metere
proposes signed admission with per-server tool allowlists~\cite{metere2026attested}, and ETDI signed, versioned tool
definitions~\cite{bhatt2025etdi}; both need servers or registries to change.

\textbf{Tool pinning.} MCP-Scan hashes tool descriptions to detect rug pulls~\cite{invariant2025mcpscan}, and
open-source proxies pin each tool's definition in a committed lockfile and block a changed listing; one also records the
package version and registry digest~\cite{mcppin,mcpseal}. All of them start the server to read its tools, and none
reports an evaluation on real attacks. Our runtime proxy (D2, Section~\ref{sec:adaptive}) follows their design.

\textbf{Agent supply chains.} Bagmar and Saraf weaponise setup instructions against coding agents and propose a
pre-install hook that checks names, sources and package age~\cite{bagmar2026setup}; it is our baseline. Kapner et al.\
report unpinned MCP declarations in 9.8\% of 2,660 agent setups~\cite{kapner2026harness}, and Shan finds that coding
assistants rarely check trust signals before they install~\cite{shan2026check}.

Table~\ref{tab:prior} compares the closest defences on the five properties MCP-Lock targets.

\begin{table}[t]
\caption{Closest defences and the properties MCP-Lock targets.\\
\textmd{P1 pins tool definitions; P2 checks before the server runs; P3 needs no change to servers or registries;
P4 checks continuity of build identity, publisher and install scripts; P5 evaluated on real attacks with a
benign-upgrade cost. \ding{51} yes, $\circ$ partly, -- no; judged from each work's own description.}}
\label{tab:prior}
\centering
\footnotesize
\setlength{\tabcolsep}{3.5pt}
\begin{tabular}{@{}lccccc@{}}
\toprule
Work & P1 & P2 & P3 & P4 & P5 \\
\midrule
Malware scanners~\cite{sejfia2022amalfi,ladisa2023crosslang} & -- & \ding{51} & \ding{51} & -- & $\circ$ \\
Signing, provenance~\cite{newman2022sigstore,slsa2023provenance,pep740} & -- & \ding{51} & -- & $\circ$ & -- \\
Package lockfiles~\cite{gamage2025lockfiles} & -- & \ding{51} & \ding{51} & $\circ$ & -- \\
Signed tool definitions~\cite{bhatt2025etdi,metere2026attested} & \ding{51} & -- & -- & -- & -- \\
MCP-Scan, pinning proxies~\cite{invariant2025mcpscan,mcppin,mcpseal} & \ding{51} & -- & \ding{51} & $\circ$ & -- \\
Pre-install hook~\cite{bagmar2026setup} & -- & \ding{51} & \ding{51} & -- & $\circ$ \\
\textbf{MCP-Lock (this work)} & \ding{51} & \ding{51} & \ding{51} & \ding{51} & \ding{51} \\
\bottomrule
\end{tabular}
\end{table}
